% --- Archon physics formalization source begin ---
% archon:physics
% archon:covers IPhO2026Problems/problem_IPhO_2026_2_C_3.lean
% archon:source-report reports/ipho_2026/problem_IPhO_2026_2_C_3.source.json
% archon:problem-id IPhO_2026_2
% archon:part-id C.3
% archon:previous-part IPhO_2026_2_C_1
% archon:previous-part-policy natural_language_prerequisite_only
% archon:previous-part IPhO_2026_2_C_2
% archon:previous-part-policy natural_language_prerequisite_only

\chapter{Physics problem IPhO\_2026\_2\_C\_3}
\label{ch:IPhO2026Problems_problem_IPhO_2026_2_C_3}

\paragraph{Problem source.}
For the half-cylindrical mirror of radius R, ray A is incident at angle theta
and its reflected line is y = m\_A*x + b\_A.  A neighboring parallel ray B is
incident at theta + Delta theta, with Delta theta much smaller than theta, and
its reflected line is y = m\_B*x + b\_B.  The envelope/intersection of neighboring
rays forms the caustic.  Use Figure 2g and its coordinate convention.

Current subquestion:
Find the limiting intersection coordinates (X\_c,Y\_c) of the neighboring reflected rays.

\paragraph{Current subquestion.}
Find the limiting intersection coordinates (X\_c,Y\_c) of the neighboring reflected rays.

\paragraph{Recorded answer/context.}
X\_c = R*sin(theta)\textasciicircum{}3; Y\_c = (R/2)*cos(theta)*(2 - cos(2*theta)).

\paragraph{Figure/image path.}
/root/proposal\_for\_physic/science-mango/ipho\_2026\_source/image/T2\_page-4.png

\paragraph{Reusable previous-part conclusions.}
\begin{itemize}
\item Source C.1. Question: Write the slope m\_A and intercept b\_A of reflected ray A in terms of theta and R. Reusable conclusions: m\_A = cot(2*theta), and b\_A = R/(2*cos(theta)). Policy: natural\_language\_prerequisite\_only; do\_not\_import\_Lean\_output
\item Source C.2. Question: Expand m\_B and b\_B to first order in Delta theta. Reusable conclusions: m\_B = cot(2*theta) - 2*csc(2*theta)\textasciicircum{}2*Delta theta; b\_B = [R/(2*cos(theta))]*(1 + tan(theta)*Delta theta), up to O(Delta theta\textasciicircum{}2). Policy: natural\_language\_prerequisite\_only; do\_not\_import\_Lean\_output
\end{itemize}

\paragraph{Formalization target.}
Redraft the existing scaffold under an explicit Physlib/PhysLean import.  The
mirror radius, point coordinates and ray intercepts are physical lengths with a
named projection to the common Figure 2g coordinate unit.  Keep the neighboring
intersection and punctured-neighborhood limit structure.  Leave proof bodies as
`by sorry`.

\begin{definition}[PhysicalLength]
  \label{decl:physics:IPhO_2026_2_C_3:PhysicalLength}
  \lean{IPhO2026Problems.IPhO2026_2_C_3.PhysicalLength}
  A physical length is a dimensionful real quantity carrying the physical
  dimension of length.
\end{definition}
\begin{proof}
  This is the defining type abbreviation for dimensionful real lengths.
\end{proof}

\begin{definition}[Figure2gLengthProjection]
  \label{decl:physics:IPhO_2026_2_C_3:Figure2gLengthProjection}
  \lean{IPhO2026Problems.IPhO2026_2_C_3.Figure2gLengthProjection}
  A Figure 2g length projection records the single unit choice used for every
  coordinate in the diagram.
\end{definition}
\begin{proof}
  This is the defining structure containing the common length-unit choice.
\end{proof}

\begin{definition}[Figure2gLengthProjection.readout]
  \label{decl:physics:IPhO_2026_2_C_3:Figure2gLengthProjection:readout}
  \lean{IPhO2026Problems.IPhO2026_2_C_3.Figure2gLengthProjection.readout}
  \uses{decl:physics:IPhO_2026_2_C_3:PhysicalLength, decl:physics:IPhO_2026_2_C_3:Figure2gLengthProjection}
  The projection reads a physical length as its real coordinate in the stored
  common unit.
\end{definition}
\begin{proof}
  Evaluate the dimensionful length in the stored unit choice and take its real
  value.
\end{proof}

\begin{definition}[Figure2gMirror]
  \label{decl:physics:IPhO_2026_2_C_3:Figure2gMirror}
  \lean{IPhO2026Problems.IPhO2026_2_C_3.Figure2gMirror}
  \uses{decl:physics:IPhO_2026_2_C_3:PhysicalLength}
  The half-cylindrical mirror in the coordinate system of Figure 2g.
\end{definition}
\begin{proof}
  This is the defining declaration; unfolding it gives the stated typed data or relation.
\end{proof}

\begin{definition}[Figure2gPoint]
  \label{decl:physics:IPhO_2026_2_C_3:Figure2gPoint}
  \lean{IPhO2026Problems.IPhO2026_2_C_3.Figure2gPoint}
  \uses{decl:physics:IPhO_2026_2_C_3:PhysicalLength}
  A point represented by its two Figure 2g coordinate readouts.
\end{definition}
\begin{proof}
  This is the defining declaration; unfolding it gives the stated typed data or relation.
\end{proof}

\begin{definition}[Figure2gMirror.OnReflectingSurface]
  \label{decl:physics:IPhO_2026_2_C_3:Figure2gMirror:OnReflectingSurface}
  \lean{IPhO2026Problems.IPhO2026_2_C_3.Figure2gMirror.OnReflectingSurface}
  \uses{decl:physics:IPhO_2026_2_C_3:Figure2gLengthProjection, decl:physics:IPhO_2026_2_C_3:Figure2gLengthProjection:readout, decl:physics:IPhO_2026_2_C_3:Figure2gMirror, decl:physics:IPhO_2026_2_C_3:Figure2gPoint}
  The reflecting upper semicircle shown in Figure 2g.
\end{definition}
\begin{proof}
  This is the defining declaration; unfolding it gives the stated typed data or relation.
\end{proof}

\begin{definition}[ReflectedRayLine]
  \label{decl:physics:IPhO_2026_2_C_3:ReflectedRayLine}
  \lean{IPhO2026Problems.IPhO2026_2_C_3.ReflectedRayLine}
  \uses{decl:physics:IPhO_2026_2_C_3:PhysicalLength}
  The supporting affine line of a reflected optical ray in Figure 2g.
\end{definition}
\begin{proof}
  This is the defining declaration; unfolding it gives the stated typed data or relation.
\end{proof}

\begin{definition}[ReflectedRayLine.Contains]
  \label{decl:physics:IPhO_2026_2_C_3:ReflectedRayLine:Contains}
  \lean{IPhO2026Problems.IPhO2026_2_C_3.ReflectedRayLine.Contains}
  \uses{decl:physics:IPhO_2026_2_C_3:Figure2gLengthProjection, decl:physics:IPhO_2026_2_C_3:Figure2gLengthProjection:readout, decl:physics:IPhO_2026_2_C_3:Figure2gPoint, decl:physics:IPhO_2026_2_C_3:ReflectedRayLine}
  A Figure 2g point lies on the supporting line of a reflected ray.
\end{definition}
\begin{proof}
  This is the defining declaration; unfolding it gives the stated typed data or relation.
\end{proof}

\begin{definition}[IsNeighboringReflectedIntersection]
  \label{decl:physics:IPhO_2026_2_C_3:IsNeighboringReflectedIntersection}
  \lean{IPhO2026Problems.IPhO2026_2_C_3.IsNeighboringReflectedIntersection}
  \uses{decl:physics:IPhO_2026_2_C_3:Figure2gLengthProjection, decl:physics:IPhO_2026_2_C_3:Figure2gPoint, decl:physics:IPhO_2026_2_C_3:ReflectedRayLine, decl:physics:IPhO_2026_2_C_3:ReflectedRayLine:Contains}
  A point is the intersection of the reflected ray at incidence angle “θ” (ray A) and the reflected ray at the neighboring angle “θ + Δθ” (ray B).
\end{definition}
\begin{proof}
  This is the defining declaration; unfolding it gives the stated typed data or relation.
\end{proof}

\begin{theorem}[Physics formalization target]
\label{thm:physics:IPhO_2026_2_C_3:target}
\lean{IPhO2026Problems.IPhO2026_2_C_3.limitingIntersectionCoordinates}
\uses{decl:physics:IPhO_2026_2_C_3:PhysicalLength, decl:physics:IPhO_2026_2_C_3:Figure2gLengthProjection, decl:physics:IPhO_2026_2_C_3:Figure2gLengthProjection:readout, decl:physics:IPhO_2026_2_C_3:Figure2gMirror, decl:physics:IPhO_2026_2_C_3:Figure2gPoint, decl:physics:IPhO_2026_2_C_3:Figure2gMirror:OnReflectingSurface, decl:physics:IPhO_2026_2_C_3:ReflectedRayLine, decl:physics:IPhO_2026_2_C_3:ReflectedRayLine:Contains, decl:physics:IPhO_2026_2_C_3:IsNeighboringReflectedIntersection}
The intersections of ray \(A\) with ray \(B\) tend, as
\(\Delta\theta\to0\) through nonzero values, to
\[
X_c=R\sin^3\theta,\qquad
Y_c=\frac R2\cos\theta\left(2-\cos(2\theta)\right),
\]
using the same named length projection for \(R,X_c,Y_c\).
\end{theorem}
\begin{proof}
For the two affine ray equations, subtract to obtain
\[
x(\Delta\theta)=-
\frac{b(\theta+\Delta\theta)-b(\theta)}
     {m(\theta+\Delta\theta)-m(\theta)}.
\]
The C.2 expansions show that the limit is \(-b'(\theta)/m'(\theta)\).
Substituting \(m'(\theta)=-2\csc^2(2\theta)\) and
\(b'(\theta)=R\tan\theta/(2\cos\theta)\), then simplifying, gives
\(X_c=R\sin^3\theta\).  Substitute this limit into
\(y=m(\theta)x+b(\theta)\) and use the double-angle identities to obtain the
formula for \(Y_c\).
\end{proof}
% --- Archon physics formalization source end ---
